\subsection{Cooperative AI Systems and Multi-Agent Moral Development}
\label{sec:5.5.1}
Agents that learn alongside other agents are learning from a social environment, and what that environment rewards is what they acquire. Multi-agent reinforcement learning makes the choice explicit: reward cooperation, adherence to a norm, and resistance to coercion directly, and a population moves toward those dispositions instead of arriving at them as a side effect of something else \autocite{leibo2017sequential}. The reward specification is a moral specification whether or not anyone writes it down as one.

Two working cases show the range. Collaborative robots share physical space with people and could acquire norms from the people they work beside, which is social learning with all of that mechanism's exposure to whatever the workplace rewards \autocite{colgate1996cobots}. Self-driving systems are validated largely in simulation, and a simulation can include forced trade-offs between competing risks, which means a decision policy is shaped by scenarios somebody chose to write \autocite{behzadan2019trolleymod}.

The limit is a closed loop. A population of agents that only ever learns from other agents has closed the loop on the question of what ethical means, and the answer it converges on is a fact about the reward function rather than about anything the affected people would recognize. Disagreement is collected from a population and resolved by a mechanism nobody outside it can move.

\runin{What opens the loop}
Three levers do most of the work of keeping it open. Open tooling lets researchers outside a lab inspect and replicate what the lab claims, in a field where claims about a model's behavior are otherwise taken on faith \autocite{wolf2020transformers}, and multi-agent research platforms are the specific case, letting norm formation and adversarial drift be studied under conditions closer to deployment than a single-agent benchmark allows \autocite{leibo2021meltingpot}. Competitions pull toward whatever they measure: a leaderboard on a fixed benchmark says nothing about ethical behavior in either direction, and a competition judged on a real problem the entrants chose themselves gives them a reason to weigh who benefits. Standing consortia work on a longer clock and are the weakest of the three, for the reasons section~\ref{sec:9.1.2} sets out. None of this works without people who are not AI researchers: neuroscience, psychology, philosophy and political science each catch failure modes the others are structurally blind to.
